
  Taking \(H\), in particular, to be the transform of \(M_S\)
  under an automorphism of \(G_S\), we conclude that \(M\) is
  characteristic in \(G\).

  Finally, if \(k\) is perfect, \(G/M\) lifts in \(G\) to a
  unipotent group, uniquely so by \SGARefLink{sga3:xvii:theorem:6:1:1:item:C}{6.1.1(C)}.
\end{enumerate}

\SGAStatement{Remark 7.2.2}{%
\begin{enumerate}
\item[\SGAInlineRefTarget{sga3:xvii:subsection:6:4:item:i}(i)] \SGARefTarget{sga3:xvii:remark:7:2:2:item:i}
  If \(k\) is not perfect, the unipotent component of
  \(G_{\bar k}\) need not be defined over \(k\), as \SGARefLink{sga3:xvii:subsection:6:4:item:a}{Example 6.4(a)}
  shows.

\item[(ii)] \SGARefTarget{sga3:xvii:remark:7:2:2:item:ii}
  In contrast to the multiplicative-type component \(M\), the
  unipotent component \(U\) is not in general characteristic in
  \(G\), whatever the characteristic of \(k\).  Of course, the
  uniqueness of the decomposition in \SGARefLink{sga3:xvii:theorem:7:2:1:item:b}{7.2.1(b)} implies that \(U\)
  is invariant under every \(k\)-automorphism of \(G\).  But if
  \(U\) and \(M\) are such that there exist a \(k\)-prescheme
  \(S\) and a nonzero \(S\)-homomorphism
  \[
    h:U_S\longrightarrow M_S
  \]
  (cf.\ \SGARefLink{sga3:xvii:remark:2-6}{2.6}), then one obtains an \(S\)-automorphism of \(G_S\),
  \[
    (u,m)\longmapsto\bigl(u,h(u)+m\bigr),
  \]
  which does not leave \(U_S\) invariant.

\item[(iii)] \SGARefTarget{sga3:xvii:remark:7:2:2:item:iii}
  If \(G\) is finite over \(k\), then under Cartier duality (\SGARefLink{sga3:xvii:remark:2-6}{2.6})
  the quotient \(G/M\) corresponds to the connected component of
  the identity in the dual \(D(G)\) of \(G\).
\end{enumerate}}
\SGARefTarget{sga3:xvii:remark:7:2:2}

\SGASubsection{7.3}{Structure of Nilpotent Affine Algebraic Groups}
\SGARefTarget{sga3:xvii:subsection:7:3}

\SGAStatement{Theorem 7.3.1}{Let \(k\) be a field and \(G\) a
nilpotent (Exposé VI\textsubscript{B}, \S8), affine, connected
algebraic \(k\)-group.  Then \(G\) has a largest subgroup \(M\) of
multiplicative type.  The group \(M\) is central and
characteristic, and \(G/M\) is a unipotent algebraic group.}
\SGARefTarget{sga3:xvii:theorem:7:3:1}

Let \(Z\) be the centre of \(G\), and let \(M\) be the largest
subgroup of multiplicative type in \(Z\), as supplied by \SGARefLink{sga3:xvii:theorem:7:2:1}{7.2.1}.
Since \(Z\) is characteristic in \(G\), and \(M\) is
characteristic in \(Z\), the group \(M\) is characteristic in
\(G\).  It remains to show that \(G/M\) is unipotent.  By
induction on the length of the ascending central series of \(G\),
this follows from the following more general lemma.

\SGAStatement{Lemma 7.3.2 (Rosenlicht)}{Let \(G\) be a connected
algebraic \(k\)-group and let \(Z\) be its centre.  Then the centre
\(Z'\) of \(G'=G/Z\) is unipotent.}
\SGARefTarget{sga3:xvii:lemma:7:3:2}

\SGAProof{Proof}
We may suppose that \(k\) is algebraically closed.  It is then
enough to show that \(Z'\) contains no subgroup isomorphic to
\(\boldsymbol\mu_\ell\), for any prime number \(\ell\), by
\SGARefLink{sga3:xvii:theorem:4:6:1:item:vi}{4.6.1(vi)}.  Let \(\boldsymbol\mu_\ell\) be such a subgroup of
\(Z'\), and let \(N\) be its inverse image in \(G\).  Since
\(\boldsymbol\mu_\ell\) is central in \(G'\), the group \(N\) is
invariant in \(G\).

\begin{enumerate}
\item[(i)] \SGARefTarget{sga3:xvii:lemma:7:3:2:proof:item:i}
  Suppose that \((\ell,p)=1\).  We can find an element
  \(x\in N(k)\) and an integer \(n\) with the following
  properties:
  \begin{enumerate}
  \item[(a)] \SGARefTarget{sga3:xvii:lemma:7:3:2:proof:item:i:a}
    \(x\) lifts a generator \(\bar x\) of
    \(\boldsymbol\mu_\ell\);
  \item[(b)] \SGARefTarget{sga3:xvii:lemma:7:3:2:proof:item:i:b}
    \(x^{\ell^n}\in Z^0(k)\).
  \end{enumerate}
  Indeed, it suffices to choose a lift of \(\bar x\) whose image
  in \((N/Z^0)(k)\) belongs to the \(\ell\)-Sylow subgroup.
  Raising to the \(\ell\)-th power in the commutative group
  \(Z^0\) is an étale morphism, so \(Z^0(k)\) is
  \(\ell\)-divisible.  Consequently, after multiplying \(x\) by
  an element of \(Z^0(k)\), we may suppose that
  \(x^{\ell^n}=0\).
